% ============================================================================
% Chapter 30: Formatting in Excel
% ============================================================================
\chapter{Formatting in Excel}
\label{ch:excel-format}

\begin{objectives}
  \item Format cells: font, colour, borders, alignment
  \item Apply number formats (currency, percentage, date)
  \item Use cell styles and conditional formatting
  \item Merge cells for headers
\end{objectives}

\bytetip[tip]{A well-formatted spreadsheet is SO much easier to read than a plain one. Let's make your data look professional!}

\section{Cell Formatting}
\label{sec:cell-format}
\index{MS Excel!formatting}

You can format cells just like text in Word --- change font, size, colour, bold, italic, and alignment. Select cells first, then use the Home tab.

\begin{shortcutbox}
\begin{tabular}{ll}
  \key{Ctrl}+\key{B} & Bold selected cells \\
  \key{Ctrl}+\key{1} & Open Format Cells dialog \\
  \key{Ctrl}+\key{Shift}+\key{\$} & Apply Currency format \\
  \key{Ctrl}+\key{Shift}+\key{\%} & Apply Percentage format \\
\end{tabular}
\end{shortcutbox}

\section{Number Formats}
\label{sec:number-formats}
\index{MS Excel!number formats}

\begin{table}[H]
\centering
\caption{Common number formats in Excel}
\label{tab:number-formats}
\renewcommand{\arraystretch}{1.3}
\begin{tabularx}{\textwidth}{l l X}
\toprule
\rowcolor{HeaderRow}
\textcolor{white}{\textbf{Format}} & \textcolor{white}{\textbf{Example}} & \textcolor{white}{\textbf{Use For}} \\
\midrule
\rowcolor{RowLight} General & 1234.5 & Default, no specific format \\
\rowcolor{RowDark} Number & 1,234.50 & Numbers with decimal places \\
\rowcolor{RowLight} Currency & \$1,234.50 & Money and prices \\
\rowcolor{RowDark} Percentage & 85\% & Ratios and proportions \\
\rowcolor{RowLight} Date & 13/04/2026 & Calendar dates \\
\rowcolor{RowDark} Text & 001234 & Text that looks like numbers (PINs, codes) \\
\bottomrule
\end{tabularx}
\end{table}

\section{Merging Cells}
\label{sec:merge}
\index{MS Excel!merge cells}

\textbf{Merge \& Centre} combines multiple cells into one large cell. It's perfect for creating headers that span multiple columns.

How to merge: Select the cells $\rightarrow$ Home tab $\rightarrow$ \textbf{Merge \& Centre} button.

\begin{warnbox}
When you merge cells, only the data in the \textbf{top-left cell} is kept. All other data in the selected range is \textbf{deleted}. Always merge \emph{before} typing your header, or make sure only one cell has data.
\end{warnbox}

\bytetip[fun]{Excel can display numbers in many ways. The number 0.85 can be shown as ``0.85,'' ``85\%,'' or even ``85/100'' (fraction format). It's the same number --- just different outfits!}

\begin{funfact}
The most popular font in Excel is \textbf{Calibri} (11pt) --- the same default font used in Word. Before 2007, the default was \textbf{Arial} (10pt). The change happened because Calibri is specifically designed for on-screen reading.
\end{funfact}

\begin{reviewqs}
  \item How do you apply bold formatting to a cell in Excel?
  \item Name four number formats available in Excel.
  \item What does Merge \& Centre do?
  \item What shortcut opens the Format Cells dialog?
  \item Why should you be careful when merging cells that contain data?
\end{reviewqs}

\begin{challenge}
Create a ``Monthly Budget'' spreadsheet with columns: Item, Cost, Category. Enter at least 8 items. Format the Cost column as \textbf{Currency}. Bold and colour the header row. Add borders to all cells. Merge cells A1:C1 for a title.
\end{challenge}
